\section{Bài tập chuyên đề}

\subsection{Củng cố}

\begin{tuluyen}
Phân loại các chất sau thành hydrocarbon, hợp chất hữu cơ không phải hydrocarbon hoặc hợp chất vô cơ:
\begin{tasks}(2)
\task \ce{CH4};
\task \ce{C2H5OH};
\task \ce{CO2};
\task \ce{C3H8};
\task \ce{CH3COOH};
\task \ce{CaCO3}.
\end{tasks}
\dapso{Hydrocarbon: \ce{CH4}, \ce{C3H8}; hữu cơ không phải hydrocarbon: \ce{C2H5OH}, \ce{CH3COOH}; vô cơ: \ce{CO2}, \ce{CaCO3}.}
\end{tuluyen}

\begin{tuluyen}
Viết công thức phân tử và công thức cấu tạo thu gọn của methane, ethane, propane, butane và ethylene.
\dapso{\ce{CH4}; \ce{C2H6}: \ce{CH3-CH3}; \ce{C3H8}: \ce{CH3-CH2-CH3}; \ce{C4H10}: hai CTCT; \ce{C2H4}: \ce{CH2=CH2}.}
\end{tuluyen}

\begin{tuluyen}
Viết phương trình hóa học của các phản ứng sau:
\begin{tasks}(2)
\task methane cháy hoàn toàn;
\task ethane cháy hoàn toàn;
\task propane cháy hoàn toàn;
\task butane cháy hoàn toàn;
\task ethylene cháy hoàn toàn;
\task ethylene tác dụng với bromine.
\end{tasks}
\dapso{Xem phần lí thuyết.}
\end{tuluyen}

\begin{tuluyen}
Đốt cháy hoàn toàn \SI{0.50}{\mole} methane. Tính số mol oxygen cần dùng và số mol carbon dioxide tạo thành.
\dapso{$n_{\ce{O2}}=\SI{1.00}{\mole}$; $n_{\ce{CO2}}=\SI{0.50}{\mole}$.}
\end{tuluyen}

\begin{tuluyen}
Đốt cháy hoàn toàn \SI{5.8}{\gram} butane. Tính khối lượng \ce{CO2} tạo thành. Cho $M_{\ce{C4H10}}=58$ g/mol.
\dapso{\SI{17.6}{\gram}.}
\begin{loigiaiBT}
\[
n_{\ce{C4H10}}=\frac{5.8}{58}=\SI{0.100}{\mole}.
\]

\[
\begin{array}{c@{\;}c@{\;}c@{\;}c@{\;}c}
\ce{2C4H10} & + & \ce{13O2} & \ce{->} & \ce{8CO2 + 10H2O}\\
0.100 & & & & 0.400
\end{array}
\]

\[
n_{\ce{CO2}}=\SI{0.400}{\mole}.
\]

\[
m_{\ce{CO2}}=0.400\times44=\SI{17.6}{\gram}.
\]

\textbf{Vậy,} khối lượng \ce{CO2} tạo thành là \SI{17.6}{\gram}.
\end{loigiaiBT}
\end{tuluyen}

\begin{tuluyen}
Đốt cháy hoàn toàn một hydrocarbon A thu được \SI{0.20}{\mole} \ce{CO2} và \SI{0.30}{\mole} \ce{H2O}. Biết lượng A đã đốt là \SI{0.10}{\mole}. Xác định công thức phân tử A.
\dapso{\ce{C2H6}.}
\begin{loigiaiBT}
\[
\begin{array}{c@{\quad\longrightarrow\quad}c@{\;+\;}c}
\ce{C_xH_y} & \ce{xCO2} & \ce{y/2H2O}\\
0.10 & 0.20 & 0.30
\end{array}
\]

\[
x=\frac{0.20}{0.10}=2,
\]
\[
\frac{y}{2}=\frac{0.30}{0.10}=3\Rightarrow y=6.
\]

\textbf{Vậy,} A có công thức phân tử \ce{C2H6}.
\end{loigiaiBT}
\end{tuluyen}

\begin{tuluyen}
Hỗn hợp X gồm methane và ethane có tổng số mol bằng \SI{0.40}{\mole}. Đốt cháy hoàn toàn X thu được \SI{0.60}{\mole} \ce{CO2}. Tính số mol mỗi khí.
\dapso{\SI{0.20}{\mole} \ce{CH4}; \SI{0.20}{\mole} \ce{C2H6}.}
\end{tuluyen}

\begin{tuluyen}
Nêu cách phân biệt methane và ethylene bằng một thuốc thử hóa học. Viết phương trình phản ứng xảy ra.
\dapso{Dùng nước bromine; ethylene làm mất màu bromine.}
\end{tuluyen}

\begin{tuluyen}
Cho \SI{0.30}{\mole} ethylene phản ứng hoàn toàn với bromine. Tính khối lượng bromine cần dùng. Cho $M_{\ce{Br2}}=160$ g/mol.
\dapso{\SI{48.0}{\gram}.}
\end{tuluyen}

\begin{tuluyen}
Một hỗn hợp LPG gồm $40\%$ propane và $60\%$ butane theo thể tích. Tính số mol oxygen cần để đốt cháy hoàn toàn \SI{2.00}{\mole} hỗn hợp.
\dapso{\SI{11.80}{\mole} \ce{O2}.}
\begin{loigiaiBT}
\[
n_{\ce{C3H8}}=2.00\times0.40=\SI{0.80}{\mole},
\]
\[
n_{\ce{C4H10}}=2.00\times0.60=\SI{1.20}{\mole}.
\]

\[
\begin{array}{c@{\;}c@{\;}c@{\;}c@{\;}c}
\ce{C3H8} & + & \ce{5O2} & \ce{->} & \ce{3CO2 + 4H2O}\\
0.80 & & 4.00 & &
\end{array}
\]
\[
\begin{array}{c@{\;}c@{\;}c@{\;}c@{\;}c}
\ce{2C4H10} & + & \ce{13O2} & \ce{->} & \ce{8CO2 + 10H2O}\\
1.20 & & 7.80 & &
\end{array}
\]

\[
n_{\ce{O2}}=4.00+7.80=\SI{11.80}{\mole}.
\]

\textbf{Vậy,} cần \SI{11.80}{\mole} oxygen.
\end{loigiaiBT}
\end{tuluyen}

\subsection{Bài tập tổng hợp}

\begin{tuluyen}
Một hydrocarbon A có $M_A=44$ g/mol. Đốt cháy hoàn toàn \SI{0.15}{\mole} A thu được \SI{0.45}{\mole} \ce{CO2}. Xác định công thức phân tử A và tính số mol \ce{H2O} tạo thành.
\dapso{A là \ce{C3H8}; $n_{\ce{H2O}}=\SI{0.60}{\mole}$.}
\end{tuluyen}

\begin{tuluyen}
Hỗn hợp X gồm methane và propane. Đốt cháy hoàn toàn \SI{0.50}{\mole} X thu được \SI{0.90}{\mole} \ce{CO2}.
\begin{enumerate}
\item Tính số mol mỗi hydrocarbon trong X.
\item Tính số mol oxygen cần dùng để đốt cháy hoàn toàn X.
\end{enumerate}
\dapso{$n_{\ce{CH4}}=\SI{0.30}{\mole}$; $n_{\ce{C3H8}}=\SI{0.20}{\mole}$; $n_{\ce{O2}}=\SI{1.60}{\mole}$.}
\begin{loigiaiBT}
Gọi số mol methane và propane lần lượt là $x,y$.

\[
x+y=0.50.
\]

\[
\begin{array}{c@{\;}c@{\;}c@{\;}c@{\;}c}
\ce{CH4} & + & \ce{2O2} & \ce{->} & \ce{CO2 + 2H2O}\\
x & & 2x & & x
\end{array}
\]
\[
\begin{array}{c@{\;}c@{\;}c@{\;}c@{\;}c}
\ce{C3H8} & + & \ce{5O2} & \ce{->} & \ce{3CO2 + 4H2O}\\
y & & 5y & & 3y
\end{array}
\]

\[
\begin{cases}
x+y=0.50,\\
x+3y=0.90
\end{cases}
\Rightarrow
\begin{cases}
x=0.30,\\
y=0.20.
\end{cases}
\]
\[
n_{\ce{O2}}=2x+5y=2\times0.30+5\times0.20=\SI{1.60}{\mole}.
\]

\textbf{Vậy,} X chứa \SI{0.30}{\mole} \ce{CH4}, \SI{0.20}{\mole} \ce{C3H8}; cần \SI{1.60}{\mole} \ce{O2} để đốt cháy hoàn toàn.
\end{loigiaiBT}
\end{tuluyen}

\begin{tuluyen}
Một hỗn hợp LPG gồm propane và butane có tổng số mol bằng \SI{0.50}{\mole}. Khi đốt cháy hoàn toàn thu được \SI{1.80}{\mole} \ce{CO2}. Tính phần trăm số mol mỗi khí trong hỗn hợp.
\dapso{$40\%$ propane; $60\%$ butane.}
\end{tuluyen}

\begin{tuluyen}
Xăng E10 chứa ethanol và hydrocarbon. Xét riêng phần hydrocarbon gồm \ce{C8H18} và \ce{C9H20} có tỉ lệ mol $3:4$.

Tính số mol \ce{CO2} sinh ra khi đốt cháy hoàn toàn \SI{0.70}{\mole} hỗn hợp hydrocarbon trên.
\source{Phỏng theo dữ kiện xăng E10 trong đề tuyển sinh chuyên Hóa Lâm Đồng 2026.}
\dapso{\SI{6.00}{\mole} \ce{CO2}.}
\end{tuluyen}

\begin{tuluyen}
Khí gas hóa lỏng chứa $40\%$ propane và $60\%$ butane theo thể tích. Tính khối lượng \ce{CO2} tạo ra khi đốt cháy hoàn toàn \SI{1.00}{\mole} hỗn hợp.
\source{Phỏng theo dữ kiện LPG trong đề tuyển sinh chuyên Hóa Lâm Đồng 2026.}
\dapso{$n_{\ce{CO2}}=\SI{3.60}{\mole}$; $m_{\ce{CO2}}=\SI{158.4}{\gram}$.}
\end{tuluyen}

\begin{tuluyen}
Một khí hydrocarbon A làm mất màu nước bromine. Đốt cháy hoàn toàn \SI{0.20}{\mole} A thu được \SI{0.40}{\mole} \ce{CO2}. Biết A là hydrocarbon đơn giản đã học. Xác định A và viết phương trình phản ứng của A với bromine.
\dapso{A là ethylene, \ce{C2H4}; \ce{CH2=CH2 + Br2 -> CH2Br-CH2Br}.}
\end{tuluyen}
